\paragraph{Bilingual questions about pictured Chinese novel passages.}
BiPaR supplies human Chinese and English questions with annotated answers in the same Chinese passage (\href{https://aclanthology.org/D19-1249/}{Jing, Xiong and Yan, 2019}, CC BY-NC 4.0). We verified answer offsets, retained at most 416 source cells and 48 answer cells, and screened the resulting passages against training and reserved evaluation material. Within each training passage we constructed a maximum assignment that exchanged questions with disjoint answers. This retained 11,451 bilingual records, producing 22,902 book presentations and 26,381 earlier reading and historical presentations in each arm. Each writer inherited the same text-pretrained Chinese transformer and learned visual input and output interface. We fitted the last four transformer blocks, final normalization and pointer, keeping the glyph correspondence fixed. The loss sums negative log probability of the next observed glyph under the mixture of inventory generation and source copying, with auxiliary supervision on the first address of a matching answer. English alignment, Chinese repetition and same-page English reassignment all used one epoch, the same seed and the same source images. Length-based batching could group the questions differently.

On each of 374 official validation passages, we selected the first two questions with disjoint answer sets in source order, then rendered each in both languages. Training and validation can share novel titles. The official test split remained unopened. A passage is complete only if the four autonomous replies stop and match recorded glyph sequences. A deterministic writer whose question representation collides across either pair cannot pass that passage. For a specified target sequence, greedy decoding is exact if and only if every action along the reference prefix, including the final stop, is correct. A reference score margin larger than twice a uniform perturbation bound preserves the selected action. We checked these statements in Lean and located the first disagreement in saved reply images, retaining every viable alternative reference at each step. Partial credit uses longest-common-subsequence F1 on physical cells. The same printed source and question identities were used for image and text serving comparisons.

Autoregressive pixel generation has previously been assessed for legibility and text generation (\href{https://aclanthology.org/2024.findings-acl.874/}{Tai et al., 2024}). Multilingual visual text encoders have also been evaluated on semantic similarity and retrieval (\href{https://aclanthology.org/2026.emnlp-main.2412/}{Yuan et al., 2026}). Here the two questions with incompatible answers provide a test of which facts the writer selects on a fixed page.
